% ---------------------------------------------------------------------------
% Summary and Discussion draft -- written to close the argument opened in
% notes/intro_draft.tex and carried through notes/methods_results_draft.tex.
% Same conventions: % CHECK / % TODO / % UNVERIFIED (see methods draft header).
%
% Each bold lead below answers one of the four "differences from de Graaff"
% promised in the introduction, in the same order, then discusses what the
% mock comparison taught us and what comes next.
% ---------------------------------------------------------------------------

\section{Summary and Discussion}
\label{sec:discussion}

We have stacked the Planck 2018 lensing convergence map on $217{,}702$,
$459{,}051$, and $1{,}282{,}969$ BOSS CMASS galaxy pairs at transverse
separations of $5$, $10$, and $20\,\hinvMpc$,
subtracted a control built from the galaxies' own single-galaxy stack, and
measured the convergence left between the pair members. The same pipeline
was run on a CMASS-like mock built from the BigMDPL simulation, which
supplies the $\Lambda$CDM expectation for every quantity we measure.

\textbf{We do not detect a filament, and the data cannot yet test the
$\Lambda$CDM prediction.} The bridge excess is consistent with zero at all
three separations, at $0.6\sigma$, $2.3\sigma$, and $0.6\sigma$.
Testing the compressed profile against zero with the full covariance gives
the same conclusion. The mock predicts a bridge of $6.32\pm0.36$,
$4.91\pm0.21$, and $2.18\pm0.12\times10^{-4}$ in the same statistic,
at or below our errors at every separation --- comfortably so at
$5\,\hinvMpc$, but by only a few per cent at $10$ and $20\,\hinvMpc$ --- so
the observations cannot distinguish the prediction from no bridge at all.
At $10\,\hinvMpc$ the measurement sits $2.3\sigma$ from zero and $1.4\sigma$
from the mock, so such preference as there is runs mildly towards the
prediction rather than away from it.
% TODO: 95 per cent upper limits as Sigma and linear mass density.
The $2.3\sigma$ excess at $10\,\hinvMpc$ is at the level of the
$1.9\sigma$ signal that \citet{deGraaff2019} found in the same data with a
different estimator, and we regard the two as the same marginal hint.

\textbf{Resolving the measurement in separation costs sensitivity that
Planck cannot afford.} The first change we made relative to
\citet{deGraaff2019} was to bin pairs by separation on a fixed comoving grid
rather than rescaling and pooling them. The mock shows why this is worth
doing: the predicted bridge falls by a factor of three from $5$ to
$20\,\hinvMpc$, so the pooled measurement mixes structures of very
different amplitude. But dividing the sample raises the error in each bin,
and at Planck noise none of the bins is individually informative. Pooling
our three bins with inverse-variance weights gives
$3.05\pm1.99\times10^{-4}$, or $1.5\sigma$, which is the number to compare
with theirs. The three samples are drawn from one catalogue over one
footprint and share galaxies, so this error, which assumes they are
independent, is a lower bound.

\textbf{The limit is the map, not the method.} Removing the control
entirely and measuring the same window on the raw band changes the error
only marginally,
% TODO: quote the exact change from the committed pipeline (the 7 per cent figure came from the fitted-control analysis).
so the error budget is set by the Planck reconstruction noise within the
bridge window and not by the uncertainty of the control. No refinement of
the template can shrink the error bar; only more pairs or a less noisy
convergence map can. Treating pixels as independent would have understated
the error by a factor of $2.2$ to $2.6$, % CHECK
which is the single largest source of over-optimism in a measurement of
this kind.

\textbf{The unit-amplitude control is the right choice, and the mock says
why.} The third change was to use a control with no fitted parameters. On
the mock, the unit-amplitude control recovers a clear bridge with the
expected decline in separation, whereas a control whose core and tail
amplitudes are fitted to the region outside the pair drives the residual to
zero at every separation. % UNVERIFIED: fitted-control code not in the repo
The pair axis is aligned with the surrounding large-scale structure on both
sides of the pair, so any region along the axis is correlated with the
bridge and a fit there absorbs it. This is the weakness of profile fitting
that we anticipated in Section~\ref{sec:intro}, now measured on a
simulation with known content. The non-physical mock pairs, with
line-of-sight separations of $20$--$120\,\hinvMpc$, show no bridge with the
unit-amplitude control, % TODO: values
so the control does not manufacture a ridge on its own.

\textbf{The control is not perfect, and the imperfection is informative.}
The observed filament map retains a residual of about $2\times10^{-3}$ at
the galaxy positions against $0.9\times10^{-3}$ in the mock. % CHECK
Galaxies with a close companion live in more massive haloes than the
average CMASS galaxy, and this selection is stronger in the data than in a
central-only mock with a soft mass threshold. Independently, the BOSS
single-galaxy profile falls below the mock at large radius. % CHECK
Both discrepancies affect the isotropic residual, which the band difference
removes to first order, but they set the size of the halo signal that leaks
toward the bridge window, and they mean the mock does not reproduce the
data's leakage exactly. Two improvements to the mock follow directly. The
HOD should include satellites, which contribute the close pairs at
$5\,\hinvMpc$ and raise the mass selection in pairs, and it should be tuned
on the measured single-galaxy lensing profile rather than on the number
density alone.

\textbf{The bridge window must clear the beam.} The window
$|X|\leq0.35\,d$ ends $0.15\,d$ from each halo centre: $0.75\,\hinvMpc$ at
$d=5$, $1.5$ at $d=10$, and $3\,\hinvMpc$ at $d=20$, against a smoothed halo
profile of $3.3\,\hinvMpc$ at half maximum. Because the central band at a
given $X$ lies closer to the halo than the off-centre band, an isotropic
halo residual leaks into the band difference with a positive sign. In an
earlier version of this analysis, excluding only $\pm1\,\hinvMpc$ around
each halo produced an apparent $2.4\sigma$ bridge that vanished once the
exclusion reached a full beam width. % CHECK: from the earlier (uncommitted) analysis; confirm the numbers or drop the sentence
At $5\,\hinvMpc$ a beam-wide exclusion leaves no window at all, which is
why we treat that bin as a consistency check on the two larger ones rather
than as an independent measurement, and why the $5\,\hinvMpc$ mock bridge
should be read as bridge plus halo leakage. % TODO: quote the mock leakage estimate (B(X) with the filament removed, or with randomised pair orientations)
The general lesson is that any pair-stacking bridge statistic needs an
exclusion of at least one smoothing scale, or an explicit leakage
calibration from a simulation with the same smoothing and the same halo
selection.

\textbf{Varying the line-of-sight cut trades completeness against
contamination without a clear optimum at Planck noise.} The fourth change
was to test the line-of-sight cut on the mock rather than adopt the
$5$--$6\,\hinvMpc$ value used previously. The redshift-space sweep gives
completeness and interloper fractions of % TODO
at cuts of $5$, $10$, $20$, and $30\,\hinvMpc$, and the mock
signal-to-noise selects $5\,\hinvMpc$ for the closest pairs and
$10\,\hinvMpc$ for the others. The gain over a uniform $5\,\hinvMpc$ cut is
% TODO
which is real in the mock but well below what would be needed to change the
observational conclusion. Adopting a single cut across all bins would make
the samples directly comparable at little cost, and we recommend it for
future work.

\textbf{Prospects.} Because the error is dominated by map noise it falls as
the number of pairs grows, so the most direct gain with Planck is a wider
separation bin, at the cost of some smearing of the bridge. The larger step
is a less noisy convergence map. The ACT DR6 reconstruction already overlaps
a large part of the CMASS footprint with several times lower noise on the
relevant scales, % CHECK: quote overlap area and noise ratio
and the Simons Observatory and CMB-S4 will extend this further; the SPT-3G
field lies largely outside the CMASS footprint. Scaling our errors by the
noise ratio, the mock bridge at $10\,\hinvMpc$ would be detected at
% TODO: forecast S/N
with the present pair sample. Combining the convergence stack with the
thermal Sunyaev--Zel'dovich stack of the same pairs, as \citet{deGraaff2019}
did, would then constrain the gas temperature of the bridge as well as its
mass. Finally, the same pairs are covered by the HSC, DES, and KiDS shape
catalogues over smaller areas, and a joint galaxy-shape and CMB-lensing
measurement would test the bridge with two independent source planes.
